% =====================================================================
% Fig 4 -- Manas-Buddhi two-stage loop with Sakshi observation
% =====================================================================
\begin{figure}[!tbp]
\centering
\resizebox{0.92\linewidth}{!}{%
\begin{tikzpicture}[
  font=\small,
  every node/.style={rounded corners=2pt},
  io/.style    ={draw=black!60,fill=gray!8,inner sep=4pt,minimum width=30mm,
                 minimum height=10mm,align=center},
  manas/.style ={draw=black!60,fill=orange!8,inner sep=4pt,minimum width=46mm,
                 minimum height=18mm,align=center},
  buddhi/.style={draw=black!60,fill=blue!8,inner sep=4pt,minimum width=46mm,
                 minimum height=18mm,align=center},
  sakshi/.style={draw=black!60,fill=violet!6,inner sep=4pt,minimum width=46mm,
                 minimum height=12mm,align=center},
  store/.style ={draw=black!60,fill=green!6,inner sep=4pt,minimum width=30mm,
                 minimum height=18mm,align=center},
  arr/.style   ={-{Latex[length=2mm]},gray!70,thick},
  obs/.style   ={-{Latex[length=2mm]},dashed,violet!70,thick},
  >=Latex
]

% --- Inputs (left) ---
\node[io] (q) at (0,0) {User turn $u_t$\\\footnotesize (NL prompt)};
\node[io,below=4mm of q] (tools) {Tool returns\\\footnotesize (files, DB rows)};

% --- Manas (top middle) ---
\node[manas,right=12mm of q] (manas) {%
   \textbf{ManasAgent}\\[2pt]
   \footnotesize \emph{select \& attend}\\
   \footnotesize 1. retrieve $K$ items by $a,Q$-slot\\
   \footnotesize 2. budget-rank by salience$\times p$\\
   \footnotesize 3. assemble window};

% --- Buddhi (right middle) ---
\node[buddhi,right=12mm of manas] (buddhi) {%
   \textbf{BuddhiAgent}\\[2pt]
   \footnotesize \emph{judge \& tag}\\
   \footnotesize 1. classify khy\=ativ\=ada\\
   \footnotesize 2. detect $a,Q$-conflict\\
   \footnotesize 3. emit answer + tags};

% --- Output (right) ---
\node[io,right=12mm of buddhi] (out) {Tagged answer $r_t$\\\footnotesize $+$ confidence};

% --- Context store (bottom middle) ---
\node[store,below=14mm of manas] (store) {Context store\\\footnotesize $\{(a,Q,R,p)\}_i$};

% --- Sakshi (bottom) ---
\node[sakshi,below=4mm of store,xshift=44mm] (sakshi) {%
   \textbf{Sak\d{s}\=\i Keeper}\\
   \footnotesize append-only per-turn log};

% --- Arrows: pipeline ---
\draw[arr] (q.east)     -- (manas.west);
\draw[arr] (tools.east) -| ([xshift=-6mm]manas.south) -- (manas.south);
\draw[arr] (manas.east) -- node[above,font=\footnotesize]{window} (buddhi.west);
\draw[arr] (buddhi.east)-- (out.west);

% --- Arrows: store reads & writes ---
\draw[arr,dashed,gray!60] (store.north -| manas.south) -- (manas.south)
  node[midway,right,font=\footnotesize]{read};
\draw[arr,dashed,gray!60] (buddhi.south) |- ([yshift=-2mm]store.east)
  node[pos=0.25,right,font=\footnotesize]{write};

% --- Arrows: Sakshi observation (passive, both stages) ---
\draw[obs] (manas.south)  to[bend right=8] (sakshi.north west);
\draw[obs] (buddhi.south) to[bend left=12] (sakshi.north);
\draw[obs] (store.south)  -- (sakshi.north -| store.south);

% --- Box ---
\begin{pgfonlayer}{background}
  \node[draw=black!30,dashed,inner sep=4mm,fit=(manas)(buddhi),
        label={[font=\footnotesize\itshape]above:Two-stage agent loop}]{};
\end{pgfonlayer}

\end{tikzpicture}}
\caption{\textbf{The Manas--Buddhi two-stage loop, observed by the Sak\d{s}\=\i Keeper.}
\textbf{Manas} (S\=a\d{m}khya: discriminative attention) reads the
context store, selects up to $K$ items by qualifier-slot, and assembles
the window under a hard token budget.
\textbf{Buddhi} (Advaita Ved\=anta: discerning judgment) consumes that
window, classifies its own answer along the six-class khy\=ativ\=ada
taxonomy (Fig.~\ref{fig:khyati}), and emits the tagged answer back to
the host.
The \textbf{Sak\d{s}\=\i Keeper} (Advaita: witness consciousness) is
strictly passive --- it never enters the response path --- and writes
one immutable JSON line per stage, including hashes of the window
actually shown to Buddhi.
This separation is what gives the system its post-hoc auditability.}
\label{fig:manas_buddhi}
\end{figure}
